% ------------------------------------------------------------
% Page 110
% ------------------------------------------------------------

\section*{LA VALLEE DU BETHUZON}

C’est la plus courte des trois vallées qui convergent vers Meyrueis, celle que je connais le mieux : à
partir de la chaussée de Roquedols jusque au dessus de la Pépinère, je l’ai parcourue, en allant à la
pêche, depuis plus de 50 ans.

Le matin, souvent de très bonne heure nous laissions nos vélos, mon père et moi, dans une remise de
Ferrussac, du Villaret ou même de Rousses, pour remonter la rivière. Et, dans l’après-midi, une fois le
panier ou la musette garnis, il n’était pas question d’entamer la descente, sans avoir passé quelques
instants au Villaret chez les Bertrand, au « château », chez les Gache, les Martin ou les Séquier :
Belton, le père, nous racontait ses chasses à la sauvagine dans le ruisseau de la Croix de Fer. A
Ferrussac arrêt chez les Jean Bertrand, la première maison, où s’affinaient, dans la cave, les
incomparables fromages de chèvre de la cousine Délie, puis chez Numa Martin dont je ne peux
oublier la haie de groseilles éclatantes, les Séquier ou les Gache !!!

Nous avions des liens de parenté, certes très lointains, mais pas oubliés et, pour les plus anciens, les
souvenirs des Couderc et des Loupiac du temps de la filature du Moulin d’Ayres, n’avaient pas encore
disparu dans les brumes du temps.

La vallée du Béthuzon et Ferrussac, en particulier, disposait d’une certaine autonomie comme
Pourcarès et Les Oubrets dans la vallée de la Brèze, Salvinsac, La Bragouse ou Gatuzieres dans celle
de la Jonte. Ces « sections » de commune possédaient, en propre, des coupes de bois : le produit de
leur vente alimentait un budget géré par la section pour du matériel ou des investissements
communautaires.

La vallée du Béthuzon, c’était aussi la vallée des trois châteaux, celle des trois moulins et bien
évidemment, celle des trois fours à pain.

\textbf{Les châteaux :}

Roquedols avec ses tours massives, sa forêt, son histoire ancienne, et, comme tout château, un
souterrain jamais retrouvé,

le « château » de Ferrussac, grosse maison de « maître » sans grand caractère, mais ayant appartenu à
un botaniste de renom, Jacques Cambessédès qui fut maire de Meyrueis

le Villaret avec ses trois fermes protégées par les deux petites tours flanquant cette bâtisse solidement
campée sur ses assises schisteuses, appelée, bien sûr, « le château du Villaret ».

\textbf{Les moulins et les fours à pain}

La vallée du Béthuzon, c’était encore la vallée des trois moulins. Chaque hameau, vivant en autarcie
partielle, pouvait ainsi disposer de sa farine de seigle pour le pain, mais aussi de mouture plus
grossière pour les bêtes. De ces trois moulins et des chaussées et béals les alimentant ne subsistent que
quelques vestiges.

Celui de Ferrussac a été en activité jusqu’à la 2\textsuperscript{o} guerre mondiale, ceux de Rousses et du Villaret se
sont arrêtés plus tôt, dans les années trente.

A Rousses on peut discerner encore, mais avec peine, 200m avant le hameau l’emplacement du tout
petit moulin en contrebas d’une cascade souvent bien maigre en été.

Au Villaret le moulin à farine avait été transformé entre les deux guerres en scierie rudimentaire ; il
fournissait également l’électricité pour les 3 ou 4 maisons, grâce à une petite turbine installée par Jules
Gache et Belton Séquier tous deux excellents bricoleurs et voisins !

\textbf{Les chantiers et les harkis}

La vallée du Béthuzon c’est elle, qui accueille le plus grand nombre
d’appelés aux Chantiers de Jeunesse en 1940-1944 à Roquedols, Rousses et la Pépinère en amont de
Rousses et 20 ans plus tard, à Ferrussac, une trentaine de Harkis avec femmes et enfants, (130
personnes environ au total).
